\RequirePackage{plautopatch}
\documentclass[uplatex,dvipdfmx,a4paper,9pt]{jsarticle}

%このようにパーセントを打つと，その行はコメントになります．メモ書きに使えますよ！
\usepackage{amsmath}
\usepackage{amssymb}
\title{格好良い数式たち}
\author{sinuture}

\begin{document}
\maketitle

閉区間\footnote{英語で interval といいます．}$I$について，以下を考えましょう． 
\begin{align}
    & here\\
    & Aspiring sincerely to an international peace
\end{align}

\begin{align}
    \frac{1}{\pi} = \frac{2\sqrt{2}}{99^2}\sum_{n=0}^{\infty} \frac{(4n)!}{(4^nn!)^4} \frac{26390n+1103}{99^4n}
\end{align}

\begin{align*}
    y = \frac{\log_e(\frac{x}{m}-sa)}{r^2}\\
    yr^2 = \log_e(\frac{x}{m}-sa)\\
    r^{yr^2} = \frac{x}{m}-sa\\
    mr^{yr^2} = x -msa\\
    me^{rry} = x - mas
\end{align*}

\begin{equation}
    \lim_{ x \rightarrow \infty } \frac{ \sin x }{x} = - e^{ i \pi } = 1
\end{equation}

調査対象は表\ref{かおｎ}に示す通りです．\\
\begin{table}[htbp]
    \centering
    \caption{調査対象}
    \label{かおｎ}
    \begin{tabular}{l|r|r}
        \hline
        おなまえ & 年齢$/$歳 & 所持金$/\mathrm{yen}$\\
        \hline
        わたし & 12 & 50000億\\
        あなた & 13 & 500\\
        きさま & 90 & 5,000\\
        \hline
    \end{tabular}
\end{table}

\pagebreak
ここからはチャバの定理の証明です．$\overrightarrow{\mathrm{AF}} = t\overrightarrow{\mathrm{AB}}, \overrightarrow{\mathrm{AE}}=s\overrightarrow{\mathrm{AC}}$とおく．$\displaystyle\frac{\mathrm{BD}}{\mathrm{DC}} = \frac{s(1-t)}{1-s}$となることを示せば，チェバの定理が示される．

BEのベクトル方程式は，
\begin{equation}
    \overrightarrow{\mathrm{AP}} = u \overrightarrow{\mathrm{AB}} + (1-u)s \overrightarrow{\mathrm{AC}}.
\end{equation}

CFのベクトル方程式は，
\begin{equation}
    \overrightarrow{\mathrm{AP}} = (1-v)t \overrightarrow{\mathrm{AB}} + v \overrightarrow{\mathrm{AC}}.
\end{equation}

よって，$\overrightarrow{\mathrm{AB}}$，$\overrightarrow{\mathrm{AC}}$の係数を連立させて$u$について解くと，$\displaystyle u = \frac{t-st}{1-st} $となるので，
\begin{equation}
    \overrightarrow{\mathrm{AP}}
\end{equation}

konn$ \displaystyle \frac{1}{x} $がはっ\ref{sone}はｈ！

\begin{align}
    p
\end{align}

\end{document}